% 04_macros.tex
% Componentes especiales, avisos, capturas de pantalla y formatos
% (Manual de usuario — ISO/IEC/IEEE 26514:2022)

% =============================================================================
% 1. FORMATO PARA NOMBRES DE BOTONES, RUTAS Y RESULTADOS
% =============================================================================

% \boton{Texto}: nomenclatura estandarizada de un control de la interfaz.
\newcommand{\boton}[1]{{\sffamily\bfseries\color{colorPrimario}[#1]}}

% \campo{Texto}: nombre de un campo de formulario o etiqueta de la interfaz.
\newcommand{\campo}[1]{{\sffamily\itshape #1}}

% \ruta{a > b > c}: ruta de navegación dentro de la aplicación.
\newcommand{\rutanav}[1]{{\sffamily\bfseries\color{grisISO}#1}}

% \comando{texto}: comandos en línea, URL cortas o valores literales.
\newcommand{\comando}[1]{{\fontfamily{pcr}\fontsize{9}{10.8}\selectfont\colorbox{grisTecnico}{\texttt{#1}}}}

% \resultado{texto}: resultado esperado de una acción (verde de éxito).
\newcommand{\resultado}[1]{\textcolor{colorSecundario}{\textbf{#1}}}

% Tipos de columna con centrado vertical (array): L izq, C centro, R der.
\newcolumntype{L}[1]{>{\raggedright\arraybackslash}m{#1}}
\newcolumntype{C}[1]{>{\centering\arraybackslash}m{#1}}
\newcolumntype{R}[1]{>{\raggedleft\arraybackslash}m{#1}}
\newcolumntype{Y}{>{\raggedright\arraybackslash}X}

% =============================================================================
% 2. CAPTURAS DE PANTALLA
% =============================================================================
% Norma de maquetación del manual: las capturas se insertan al 0.9 del ancho
% de texto y, mediante un tope de altura, se garantiza que NO entren más de
% dos imágenes por página (evita desbordes y páginas con grandes huecos).

% Contador de capturas en la página actual.
\newcounter{capturasEnPagina}
\setcounter{capturasEnPagina}{0}

% \captura{ruta_imagen}{leyenda descriptiva}
% Inserta una captura centrada (ancho 0.9\textwidth, alto máx. 0.42\textheight)
% con su leyenda numerada. Tras la segunda captura de la página fuerza un salto
% para respetar el máximo de dos imágenes por página.
\newcommand{\captura}[2]{%
  \par\vspace{0.1cm}%
  \begin{center}%
    \includegraphics[width=0.9\textwidth,height=0.42\textheight,keepaspectratio]{#1}%
    \captionof{figure}{#2}%
  \end{center}%
  \vspace{0.1cm}%
  \stepcounter{capturasEnPagina}%
  \ifnum\value{capturasEnPagina}>1\relax
     \setcounter{capturasEnPagina}{0}%
     \clearpage
  \fi%
}

% \reiniciarPagina: reinicia el contador e inicia página nueva. Útil al cambiar
% de sección para no arrastrar una captura suelta a un bloque distinto.
\newcommand{\reiniciarPagina}{\setcounter{capturasEnPagina}{0}\clearpage}

% \responsivePar{img1}{cap1}{img2}{cap2}: dos capturas verticales (móvil) una
% al lado de la otra. Pensado para el apartado de diseño adaptable (responsive).
\newcommand{\responsivePar}[4]{%
  \par\vspace{0.3cm}%
  \noindent\begin{center}%
  \begin{minipage}[t]{0.46\textwidth}\centering
    \includegraphics[width=\linewidth,height=0.62\textheight,keepaspectratio]{#1}%
    \captionof{figure}{#2}%
  \end{minipage}\hfill
  \begin{minipage}[t]{0.46\textwidth}\centering
    \includegraphics[width=\linewidth,height=0.62\textheight,keepaspectratio]{#3}%
    \captionof{figure}{#4}%
  \end{minipage}%
  \end{center}\par\vspace{0.3cm}\clearpage}

% \responsiveUna{img}{cap}: una sola captura vertical (móvil) centrada.
\newcommand{\responsiveUna}[2]{%
  \par\vspace{0.3cm}\begin{center}%
    \includegraphics[width=0.42\textwidth,height=0.62\textheight,keepaspectratio]{#1}%
    \captionof{figure}{#2}%
  \end{center}\par\vspace{0.3cm}}

% =============================================================================
% 3. COMPONENTES ESPECIALES DE USABILIDAD (avisos ISO)
% =============================================================================

% Nota: consejos de productividad (fondo azul claro, icono info).
\newtcolorbox{AvisoNota}{
    colback=bgNota, colframe=bgNota!80!black, boxrule=0.5pt, leftrule=4pt,
    arc=2pt, title={\faInfoCircle\ \textbf{Nota}}, coltitle=black,
    fonttitle=\sffamily, attach title to upper=\quad
}

% Precaución: evitar errores operacionales (borde amarillo, texto negrita).
\newtcolorbox{AvisoPrecaucion}{
    colback=white, colframe=colorAcento, boxrule=1.5pt, arc=2pt,
    fontupper=\bfseries, title={\faExclamationTriangle\ \textbf{Precaución}},
    coltitle=black, fonttitle=\sffamily, attach title to upper=\par
}

% Advertencia: prevenir pérdida de datos (fondo rojo suave, icono peligro).
\newtcolorbox{AvisoAdvertencia}{
    colback=bgAdvertencia, colframe=red!70!black, boxrule=0.5pt, leftrule=4pt,
    arc=2pt, title={\faRadiation\ \textbf{Advertencia}}, coltitle=black,
    fonttitle=\sffamily, attach title to upper=\par
}

% Nota Técnica: dependencias o efectos del sistema (borde azul, fondo blanco).
\newtcolorbox{NotaTecnica}{
    colback=white, colframe=bordeNotaTecnica, boxrule=1.2pt, arc=2pt,
    title={\faInfoCircle\ \textbf{Nota Técnica}}, coltitle=bordeNotaTecnica,
    fonttitle=\sffamily, attach title to upper=\par\vspace{0.2em}
}

% =============================================================================
% 4. LISTAS DE VERIFICACIÓN
% =============================================================================
\newlist{checklist}{itemize}{1}
\setlist[checklist]{label=$\square$, leftmargin=*, noitemsep}

% =============================================================================
% 5. TABLA GENÉRICA DE APOYO (matriz de roles, errores, FAQ, glosario)
% =============================================================================
% Nota: se usa tabular (no tabularx) porque tabularx no puede envolverse dentro
% de un \newenvironment (realiza una relectura del cuerpo buscando su \end).
\newenvironment{TablaInfo}[2]{%
    \renewcommand{\arraystretch}{1.4}%
    \vspace{0.3cm}\noindent
    \fontsize{9.5}{12}\selectfont\sffamily
    \begin{tabular}{|>{\bfseries\color{colorPrimario}\raggedright\arraybackslash}p{3.8cm}|>{\raggedright\arraybackslash}p{11.2cm}|}
    \hline\rowcolor{headerTabla}
    {\sffamily\bfseries #1} & {\sffamily\bfseries #2} \\ \hline
}{%
    \end{tabular}\vspace{0.3cm}%
}
